\section{Limits and open problems}\label{sec:open}

We collect the main limitations of the results and the open problems they leave. The hardest of them is the first.

\paragraph{1. Affordable soundness when the learner must invent language.} Every soundness guarantee here is relative to \emph{realizability}: the target calculus, or the latent formalization, must lie in the hypothesis class. History did not respect any fixed class. It invented uniform convergence, ideals, $\varepsilon$--$\delta$ and the fundamental group as proof-generated concepts.
\begin{itemize}
\item A universal, language-inventing class restores realizability, but sound verification then costs on the order of $2^{K(h^*)}$ escalations (\cref{thm:caution:unstructured,thm:informal:escalation}).
\item What is needed is a class that is closed under the definitional extensions mathematics actually made, and that still has small escalation dimension. Equivalently, we need a theory of when concept invention, for example by anti-unification over definitions or by naming the hidden lemma found by descent, is cheap.
\end{itemize}

\paragraph{2. Computational cost.} Our bounds count samples, escalations, detections and refutation-oracle calls, not computation.
\begin{itemize}
\item Membership in the version-space intersection for untagged unions is likely NP-hard.
\item Consistency with negative bags is a multiple-instance problem. Its hardness is the formal shadow of the obfuscated-arguments problem in debate.
\item Refutation search is only semi-decidable. Our depth-relative soundness (\cref{thm:twotier:depth}) is the honest form of what can be had.
\end{itemize}
A complexity-theoretic version of the two-tier theorem would be valuable. It would identify which audits are polynomial and for which classes.

\paragraph{3. The reading layer.} In physics the irreducible judgment is the \emph{reading}: which admissible worlds the problem text determines (\cref{thm:physics:judgment}).
\begin{itemize}
\item Determinacy refutes only over-wide readings (\cref{prop:physics:gricean}).
\item What positive data identify the intended admissible class? We conjecture that official solutions together with marking schemes form a tell-tale within a class of ``minimal completions consistent with the intended imports''.
\item Learning readings soundly is, we think, the central open problem for principled checking beyond mathematics.
\end{itemize}

\paragraph{4. Sharp combinatorics.} Several bounds have gaps:
\begin{itemize}
\item The binomial conjecture for $k$-unions of intersection-closed classes (\cref{conj:caution:binomial}) would close the gap for untagged schema learning.
\item The bag-price conjecture $M_{\mathrm{bag}}^{(r)}\le r\cdot\mathrm{Ldim}$ (\cref{conj:informal:bagprice}) would give the exact exchange rate between paradoxes and counterexample objects.
\item The query complexity of sound blame, i.e.\ computing the union of the minimal conflicts with a conflict oracle, is open.
\end{itemize}

\paragraph{5. Coherence beyond Post-completeness.}
\begin{itemize}
\item For which first-order or modal targets does every invalid uniform rule have a short coherence witness in a natural designated context?
\item Is there a hierarchy of coherence-witness complexity?
\item Which logics and theories are \emph{coherence-pinned}, meaning that the residue is empty for natural families of designated contexts?
\item For arithmetic: is there a principled learner over reflection progressions that converges to ``the right amount of reflection'' under a natural feedback notion?
\end{itemize}

\paragraph{6. Simplicity and normativity.}
\begin{itemize}
\item For which classes of simple models does the product-model frontier kink, beyond parities?
\item Is there a principled, data-free choice of the simplicity coefficient, for example one tied to the cheapest postulate-violating error?
\item Which assumptions on a performance model, such as errors being time-bounded shortcuts, make joint description length identify competence despite the known non-identifiability of preferences from behaviour?
\end{itemize}

\paragraph{7. Probabilistic and contextual coherence.}
\begin{itemize}
\item Find a finitely presented proof system for valid counting sequents that is complete for de Finetti coherence.
\item Prove a probabilistic local-to-global theorem for context systems on tree covers.
\item Decide whether contextual frustration between idealized contexts should be penalized or tolerated.
\end{itemize}

\paragraph{8. Scale.}
\begin{itemize}
\item Our experiments are illustrations in small symbolic domains with realizable targets, cheap world feedback and few seeds.
\item Whether learned calculi with cautious acceptance remain affordable at the scale of real mathematical corpora (for example Lean or Metamath step data with injected fallacies at measured human rates) is an empirical question we have not addressed.
\item We have deliberately not built a strong prover. The motivating notes regard publishing insightful research on AI theorem provers as gravely dangerous. The results here concern checking, and we make no recommendations for prover development.
\end{itemize}
